% Originally: content/week-07.tex, \section{Solutions}

\section{Solutions}
\label{sec:imft_solutions}

% Transition: Claude Opus 5 (High)
Not every solution in this chapter is collected here. \Cref{ex:implicit_system_yz_of_x}
(\cpageref{ex:implicit_system_yz_of_x}) and \cref{ex:implicit_xyz_equation}
(\cpageref{ex:implicit_xyz_equation}) are solved where they appear.

\begin{exercisesolution}[Solution to \cref{ex:7.5}]
% Generator: Claude Sonnet 5 (Medium)
Write the system as $F(x,y,z,u,v) = 0$ for $F : \mathbb{R}^5 \to \mathbb{R}^2$,
\[F(x,y,z,u,v) := \begin{pmatrix} xy^5+yu^5+zv^5-1 \\ x^5y+y^5u+z^5v-1 \end{pmatrix}.\]
The base point does lie on the level set, which the theorem requires before anything else:
\[F(0,1,1,1,0) = \begin{pmatrix} 0 + 1 + 0 - 1 \\ 0 + 1 + 0 - 1\end{pmatrix}
  = \begin{pmatrix} 0 \\ 0\end{pmatrix}.\]
The differential is
\[\Jac F(x,y,z,u,v) = \begin{pmatrix} y^5 & 5xy^4+u^5 & v^5 & 5yu^4 & 5zv^4 \\ 5x^4y & x^5+5y^4u & 5z^4v & y^5 & z^5 \end{pmatrix},\]
so at the base point,
\[\Jac F(0,1,1,1,0) = \begin{pmatrix} 1 & 1 & 0 & 5 & 0 \\ 0 & 5 & 0 & 1 & 1 \end{pmatrix}.\]
By the implicit function theorem applied with $(x,y,z)$ as the free variables and $(u,v)$ as the
variables to solve for, we need the submatrix of partial derivatives with respect to $u,v$ to be
invertible; this is
\[D_{(u,v)}F(0,1,1,1,0) = \begin{pmatrix} 5 & 0 \\ 1 & 1\end{pmatrix}, \qquad \det = 5 \neq 0,\]
so the implicit function theorem applies, and the system is indeed solvable for $u,v$ as
$C^\infty$ functions of $(x,y,z)$ near $(0,1,1)$. Writing $G(x,y,z) := (u(x,y,z),v(x,y,z))$, the
theorem gives
\begin{align*}
\Jac G(0,1,1)
  &= -\begin{pmatrix}5&0\\1&1\end{pmatrix}^{-1}\begin{pmatrix}1&1&0\\0&5&0\end{pmatrix} \\
  &= -\frac15\begin{pmatrix}1&0\\-1&5\end{pmatrix}\begin{pmatrix}1&1&0\\0&5&0\end{pmatrix} \\
  &= -\frac15\begin{pmatrix}1&1&0\\-1&24&0\end{pmatrix}.
\end{align*}
So, $D_{(0,1,1)}u = -\tfrac15\left(\begin{smallmatrix}1&1&0\end{smallmatrix}\right)$ and
$D_{(0,1,1)}v = -\tfrac15\left(\begin{smallmatrix}-1&24&0\end{smallmatrix}\right)$.
\end{exercisesolution}

\begin{exercisesolution}[Proof of \cref{ex:implicit_function_arctan}]
% Generator: Gemini 3.6 Flash (High)
\begin{enumerate}[label=\textbf{(\alph*)}]
  \item Define $F(x,y,z) := x^2 \arctan(z) + z - y$ for $(x,y,z) \in \mathbb{R}^3$. $F$ is smooth on $\mathbb{R}^3$.
For any fixed $(x,y) \in \mathbb{R}^2$, consider $h(z) := F(x,y,z) = x^2 \arctan(z) + z - y$. Its derivative with respect to $z$ is:
\[\frac{\partial F}{\partial z}(x,y,z) = \frac{x^2}{1+z^2} + 1 \ge 1 > 0 \quad \text{for all } z \in \mathbb{R}.\]
Since $\partial_z F > 0$, $h(z)$ is strictly increasing on $\mathbb{R}$. As $z \to \pm\infty$, $h(z) \to \pm\infty$. By the Intermediate Value Theorem, there exists a unique $z \in \mathbb{R}$ such that $F(x,y,z) = 0$.
Thus, for every $(x,y)$, there is a unique value $f(x,y) := z$. By the Implicit Function Theorem (\cref{thm:implicit_function_theorem}), since $\partial_z F \neq 0$ everywhere, $f$ is continuously differentiable on $\mathbb{R}^2$.

  \item Setting $y = 0$ in the defining equation:
\[0 = x^2 \arctan(f(x,0)) + f(x,0).\]
Since $z = 0$ is a solution to $x^2 \arctan(z) + z = 0$ and the solution $z$ is unique, we conclude $f(x,0) = 0$ for all $x \in \mathbb{R}$.

To find $\nabla f(x,0)$, differentiate the identity $x^2 \arctan(f(x,y)) + f(x,y) - y = 0$ implicitly
and then evaluate at $(x,0)$, where $f = 0$ and hence $\arctan(f) = 0$ and $1/(1+f^2) = 1$:
% Correction: Claude Opus 5 (High) --- "w.r.t." expanded per the shorthand rule, ^T -> \transp, and
% the chains of \implies rewritten as prose.
\begin{itemize}
  \item \textbf{With respect to $x$:} the identity gives $2x \arctan(f) + x^2 \frac{\partial_1 f}{1+f^2} + \partial_1 f = 0$. At $(x,0)$ the first term drops out, leaving $(x^2+1)\,\partial_1 f(x,0) = 0$, and since $x^2 + 1 > 0$ we conclude $\partial_1 f(x,0) = 0$.
  \item \textbf{With respect to $y$:} the identity gives $x^2 \frac{\partial_2 f}{1+f^2} + \partial_2 f - 1 = 0$. At $(x,0)$ this becomes $(x^2+1)\,\partial_2 f(x,0) = 1$, so $\partial_2 f(x,0) = \frac{1}{x^2+1}$.
\end{itemize}
Therefore $\nabla f(x,0) = \left(0, \frac{1}{x^2+1}\right)\transp$.
\end{enumerate}
\begin{remark}
Part \textbf{(a)} is stronger than what the implicit function theorem alone delivers. The theorem is
local: it produces an $f$ near a given solution point. What upgrades that to a single $f$ defined on
all of $\mathbb{R}^2$ is the strict monotonicity of $z \mapsto F(x,y,z)$ together with its unbounded
range, which pins down exactly one root for every $(x,y)$. The theorem then supplies the regularity.
\end{remark}
\end{exercisesolution}

\begin{exercisesolution}[Solution of \cref{ex:kobel_conic_implicit_maximal_interval}]
% Generator: Claude Opus 5 (High)
Write $F(x,y) := 2x^2 - 4xy + y^2 - 3x + 4y$, a polynomial and hence $C^\infty$ on
$\mathbb{R}^2$.

\textbf{(a)} First, $(1,1)$ lies on the set: $F(1,1) = 2 - 4 + 1 - 3 + 4 = 0$. Next,
\[\partial_yF(x,y) = -4x + 2y + 4, \qquad \partial_yF(1,1) = -4+2+4 = 2 \neq 0.\]
By \cref{thm:implicit_function_theorem} there are open intervals $I \ni 1$ and $J \ni 1$ and a
$C^\infty$ function $g : I \to J$ with $g(1)=1$ whose graph is exactly
$\{F=0\} \cap (I\times J)$.

\textbf{(b)} Differentiating $F(x,g(x)) = 0$ and solving, or quoting the formula in
\cref{thm:implicit_function_theorem} directly,
\[g'(x) = -\frac{\partial_xF(x,g(x))}{\partial_yF(x,g(x))},
  \qquad \partial_xF(x,y) = 4x - 4y - 3 .\]
At the point in question $\partial_xF(1,1) = 4-4-3 = -3$ and $\partial_yF(1,1) = 2$, so
\[g'(1) = -\frac{-3}{2} = \boxed{\ \tfrac32\ }.\]

\textbf{(c)} The hint is the point of the part: $F$ is quadratic in $y$, so the equation can simply
be solved. Reading $F=0$ as $y^2 + (4-4x)y + (2x^2-3x) = 0$ and applying the quadratic formula,
\[y = (2x-2) \pm \sqrt{(2x-2)^2 - (2x^2-3x)} = (2x-2) \pm \sqrt{2x^2-5x+4}.\]
The radicand $2x^2-5x+4$ has discriminant $25 - 32 = -7 < 0$ and positive leading coefficient, so
it is \emph{strictly positive for every real $x$}. Both branches are therefore defined on all of
$\mathbb{R}$, and they never meet. At $x=1$ they take the values $0 \pm 1$, so the branch through
$(1,1)$ is the one with the plus sign:
\[g(x) = 2x - 2 + \sqrt{2x^2-5x+4}, \qquad \text{maximal domain } \boxed{\ \mathbb{R}\ }.\]
As a check on \textbf{(b)}, $g'(x) = 2 + \dfrac{4x-5}{2\sqrt{2x^2-5x+4}}$ gives
$g'(1) = 2 - \tfrac12 = \tfrac32$, as found above.
\begin{remark}[Local theorem, global answer]
\label{rem:implicit_local_versus_global}
The two halves of this exercise are worth holding side by side, because they answer different
questions and it is easy to think the first has answered the second.

\Cref{thm:implicit_function_theorem} is \emph{local}, and deliberately so: it produces some
interval around $x=1$ and says nothing whatever about how far that interval extends. Part
\textbf{(c)} shows the honest answer here is \qt{all of $\mathbb{R}$}, and no amount of squeezing
the theorem would have produced it. What produced it was solving the equation.

The reverse is also worth noting. The condition $\partial_yF \neq 0$ fails along the line
$y = 2x-2$, which is precisely where the two branches would meet --- and they never do, because
the radicand never vanishes. So the hypothesis of the theorem happens to hold along the whole
branch, which is why the branch stays a graph forever. When a conic \emph{does} have a vertical
tangent, that is exactly where $\partial_yF$ vanishes and the local statement is all there is.
\end{remark}
\end{exercisesolution}


